% topic: algorithms
\question \textbf{ระบบเลขฐานสอง (Binary)} คือระบบตัวเลขที่ประกอบด้วยตัวเลขเพียง 2 ตัว ได้แก่ \textbf{0} และ \textbf{1} เท่านั้น โดยแต่ละหลัก (บิต) แทนค่าประจำหลักเป็นเลขยกกำลังของ 2 โดยนับจากขวามาซ้าย เริ่มต้นที่ $2^0 = 1$ ดังตาราง

\begin{center}
\begin{tabular}{|c|*{8}{c|}}
\hline
\textbf{เลขยกกำลัง} & $2^7$ & $2^6$ & $2^5$ & $2^4$ & $2^3$ & $2^2$ & $2^1$ & $2^0$ \\
\hline
\textbf{ค่าประจำหลัก} & 128 & 64 & 32 & 16 & 8 & 4 & 2 & 1 \\
\hline
\end{tabular}
\end{center}

ตัวอย่างเช่น เลขฐานสอง \texttt{1011} แปลงเป็นฐานสิบได้ดังนี้ \\
\quad $2^3$\quad\;\; $2^2$\quad\;\; $2^1$\quad\;\; $2^0$ \\
\quad \; 1 \qquad 0 \qquad 1 \qquad 1 \;\; $= (1 \times 8) + (0 \times 4) + (1 \times 2) + (1 \times 1) = 8 + 0 + 2 + 1 = 11$

\textbf{ขั้นตอนวิธี} สำหรับแปลงเลขฐานสองเป็นเลขฐานสิบ (ประมวลผลจากซ้ายไปขวา) มีดังนี้

\begin{pseudocode}
\textbf{ข้อมูลนำเข้า:} สตริง $S$ ซึ่งประกอบด้วยเลขฐานสอง (อักขระ '0' และ '1' เท่านั้น) \\
\textbf{ข้อมูลส่งออก:} ค่าฐานสิบที่เป็นผลลัพธ์
\newline
1) กำหนดให้ $result = 0$ และ $i = 0$ (ตำแหน่งแรกของสตริง จากซ้ายสุด)
\newline
2) \textbf{ทำซ้ำ}ตราบใดที่ $i <$ ความยาวของสตริง $S$:
\newline
.\quad 2.1) $result = result \times 2$
\newline
.\quad 2.2) ถ้า $S[i] =$ \texttt{'1'}:
\newline
.\quad.\qquad $result = result + 1$
\newline
.\quad 2.3) $i = i + 1$
\newline
3) แสดงค่า $result$
\end{pseudocode}

\textbf{ตัวอย่างการทำงาน:} ถ้า $S = \text{"1101"}$

\begin{center}
\begin{tabular}{|c|c|c|c|}
\hline
\textbf{รอบที่} ($i$) & $S[i]$ & \textbf{การคำนวณ} & \textbf{ค่า result} \\
\hline
0 & \texttt{1} & $0 \times 2 = 0$ แล้ว $+ 1$ (เพราะเจอ \texttt{1}) & 1 \\
1 & \texttt{1} & $1 \times 2 = 2$ แล้ว $+ 1$ (เพราะเจอ \texttt{1}) & 3 \\
2 & \texttt{0} & $3 \times 2 = 6$ (เจอ \texttt{0} จึงไม่บวก) & 6 \\
3 & \texttt{1} & $6 \times 2 = 12$ แล้ว $+ 1$ (เพราะเจอ \texttt{1}) & 13 \\
\hline
\end{tabular}
\end{center}

ผลลัพธ์ที่ได้คือ \textbf{13}

\textbf{คำถาม:} ถ้า $S = \text{"10110101101"}$ เมื่อดำเนินการตามขั้นตอนวิธีข้างต้น จะได้ผลลัพธ์เป็นเลขฐานสิบที่มีค่าเท่าใด
\newline\newline
ตอบ ในกระดาษคำตอบ \newline
% answer: 1453
% solution: ไล่ตามขั้นตอนวิธีทีละบิต (11 บิต)
%   i=0:  0×2=0,  +1 = 1
%   i=1:  1×2=2,  +0 = 2
%   i=2:  2×2=4,  +1 = 5
%   i=3:  5×2=10, +1 = 11
%   i=4: 11×2=22, +0 = 22
%   i=5: 22×2=44, +1 = 45
%   i=6: 45×2=90, +0 = 90
%   i=7: 90×2=180,+1 = 181
%   i=8:181×2=362,+1 = 363
%   i=9:363×2=726,+0 = 726
%   i=10:726×2=1452,+1=1453
%   วิธีลัด (สำหรับครู): 1024+256+128+32+8+4+1 = 1453
